\documentclass[11pt,a4paper]{article}

% --- Codificación y idioma ---
\usepackage[utf8]{inputenc}
\usepackage[T1]{fontenc}
\usepackage[spanish]{babel}

% --- Paquetes generales ---
\usepackage{geometry}
\geometry{top=2.5cm, bottom=2.5cm, left=2.5cm, right=2.5cm}
\usepackage{graphicx}
\usepackage{booktabs}
\usepackage{amsmath}
\usepackage{xcolor}
\usepackage{listings}
\usepackage{hyperref}

% --- Estilo para código Arduino ---
\definecolor{codegreen}{rgb}{0,0.6,0}
\definecolor{codegray}{rgb}{0.5,0.5,0.5}
\definecolor{codepurple}{rgb}{0.58,0.0,0.82}
\definecolor{backcolour}{rgb}{0.95,0.95,0.92}

\lstdefinestyle{arduinostyle}{
  backgroundcolor=\color{backcolour},
  commentstyle=\color{codegreen},
  keywordstyle=\color{magenta},
  numberstyle=\tiny\color{codegray},
  stringstyle=\color{codepurple},
  basicstyle=\ttfamily\footnotesize,
  breakatwhitespace=false,
  breaklines=true,
  captionpos=b,
  keepspaces=true,
  numbers=left,
  numbersep=5pt,
  showspaces=false,
  showstringspaces=false,
  showtabs=false,
  tabsize=2
}
\lstset{style=arduinostyle}

% --- Título ---
\title{\textbf{Informe de Avances: Torre Hidropónica Automatizada}\\
\large Proyecto Inicial (IWG400) -- Universidad Técnica Federico Santa María}
\author{\textbf{Grupo de Trabajo:} Ingeniería Civil Electrónica}
\date{\today}

\begin{document}

\maketitle

\begin{abstract}
Este documento sintetiza el estado de avance del proyecto de automatización para agricultura familiar en el Valle del Aconcagua y zonas de escasez hídrica. Se presenta el diagnóstico territorial obtenido mediante entrevistas a usuarios y expertos, la formulación del desafío de innovación y el desarrollo técnico del prototipo 1.0 basado en una Torre Hidropónica con Arduino UNO R4 WiFi.
\end{abstract}

\section{Etapa 1: Empatizar y Contexto Territorial}
La Región de Valparaíso y el Norte Chico enfrentan una profunda crisis por la megasequía, caracterizada por una severa disparidad hídrica entre la agroindustria de exportación y la agricultura familiar campesina. A través de entrevistas semiestructuradas a actores del territorio (Panquehue, San Felipe y Tilama), se identificó que las principales barreras de adopción tecnológica son el alto costo inicial, la falta de electricidad estable en zonas rurales alejadas y la complejidad de mantenimiento.

\section{Etapa 2: Definir el Desafío}
A partir de la información recopilada, el equipo definió el siguiente desafío de innovación:
\begin{quote}
\textbf{¿Cómo podríamos} diseñar un sistema de control de cultivo e irrigación automatizado, robusto y de bajo costo, que funcione de manera sencilla e independiente para proteger la producción de los pequeños agricultores rurales frente a la escasez hídrica?
\end{quote}

\section{Etapa 3: Idear y Propuesta de Solución}
Tras explorar múltiples alternativas (riego por goteo, nebulización, malla sombra), se seleccionó el desarrollo de una \textbf{Torre Hidropónica Automatizada}. Esta solución maximiza la eficiencia hídrica al recircular el agua en un circuito cerrado sin pérdida por infiltración en suelo, permitiendo además el control activo del microclima (temperatura y sombreado).

\section{Etapa 4: Prototipar (Desarrollo Técnico 1.0)}
El prototipo funcional se implementó sobre la plataforma \textbf{Arduino UNO R4 WiFi}. La arquitectura incluye:
\begin{itemize}
  \item \textbf{Sensores de Temperatura:} Dos sensores impermeables DS18B20 para monitoreo continuo de la temperatura ambiental del aire (Pin 2) y la temperatura del estanque de agua (Pin 3).
  \item \textbf{Control de Malla de Sombra:} Servo motor de rotación continua (Pin 9) que despliega la malla UV cuando la temperatura del aire supera los $27\,^\circ\mathrm{C}$ y la abre al descender a $24\,^\circ\mathrm{C}$ (control con histéresis).
  \item \textbf{Ciclo de Riego Eficiente:} Control de bomba sumergible mediante relé de 5\,V (Pin 7) en ciclos no bloqueantes ($15\,\mathrm{s}$ activada / $60\,\mathrm{s}$ desactivada).
  \item \textbf{Interfaz de Usuario y Telemetría:} Visualización local en pantalla OLED I2C de $0{,}96''$ (SSD1306) y comandos de calibración/monitoreo por puerto serie a 115200 baudios.
\end{itemize}

\subsection{Lógica de Control en Arduino}
Se estructuró el código mediante estados finitos y tiempos no bloqueantes (\texttt{millis()}) para garantizar lecturas continuas sin pausar la ejecución del sistema:

\begin{lstlisting}[language=C++]
// Fragmento del ciclo principal de control (Arduino UNO R4 WiFi)
void loop() {
  leerSensores();       // Mide temperatura de aire y agua
  actualizarServo();    // Controla apertura/cierre de malla UV
  controlBomba();       // Maneja el ciclo ON/OFF de la bomba
  actualizarPantalla(); // Refresca pantalla OLED SSD1306
  leerSerial();         // Recibe comandos de calibracion manual
}
\end{lstlisting}

\section{Etapa 5: Plan de Testeo y Próximos Pasos}
El plan de testeo evaluará la autonomía del sistema frente a cortes de energía (integrando respaldo por panel solar), la claridad de las alertas en la pantalla OLED para usuarios rurales y la estabilidad térmica del estanque para asegurar la salud radicular del cultivo.

\end{document}
